% ==============================================================================
% Course Syllabus: MSQF 536 Machine Learning Techniques
% Formatted to fit on a single, clean official syllabus page
% ==============================================================================

\newpage
\thispagestyle{plain}

\begin{center}
    {\large\textbf{MSQF 536: MACHINE LEARNING TECHNIQUES}} \hfill {\large\textbf{CREDITS: 4}}
\end{center}
\vspace{-0.6em}
\hrule height 0.8pt
\vspace{0.6em}

\noindent\textbf{Course Objectives:}
\begin{itemize}[leftmargin=1.5em, label=$\bullet$, noitemsep, topsep=1pt, parsep=1pt]
    \item \textit{To provide knowledge of machine learning techniques.}
    \item \textit{This course familiarizes in supervised and unsupervised learning.}
\end{itemize}

\vspace{0.4em}
\noindent\textbf{Course Outcomes (COs):}
\begin{description}[leftmargin=1.4cm, labelsep=0.4em, itemsep=1pt, topsep=1pt, parsep=0pt, font=\normalfont\textbf]
    \item[CO1:] Learn the basic concepts of machine learning
    \item[CO2:] Understand the various supervised learning methods
    \item[CO3:] Knowledge of unsupervised learning algorithms
    \item[CO4:] Application of dimensionality reduction methods
    \item[CO5:] Study the fundamental concepts of Artificial Neural Network
\end{description}

\vspace{0.4em}
\noindent\textbf{Unit I: Data Pre-processing}\\
\small
Data Pre-processing --- Different types of variables --- Preparing training and test datasets --- Missing Value Imputation --- Imbalance in categorical data.
\normalsize

\vspace{0.3em}
\noindent\textbf{Unit II: Supervised Learning}\\
\small
Supervised Learning --- Framework of Regression --- Model --- Loss Function --- Logistic Regression --- K Nearest Neighbour (KNN) --- Naïve Bayes --- Tree based models --- Decision Tree --- Random Forest --- Support Vector Machine.
\normalsize

\vspace{0.3em}
\noindent\textbf{Unit III: Unsupervised Learning}\\
\small
Unsupervised Learning --- Cluster Analysis --- Similarity and Dissimilarity measures --- Proximity measures --- Types of clustering --- Hierarchical clustering: Agglomerative methods and Divisive methods --- Non-hierarchical clustering: k-means clustering algorithm --- Density based clustering --- DBSCAN.
\normalsize

\vspace{0.3em}
\noindent\textbf{Unit IV: Dimensionality Reduction \& Ensemble Methods}\\
\small
Dimensionality Reduction --- Principal Component Analysis --- Singular Value Decomposition (SVD) --- Reinforcement Learning --- Recurrent Reinforcement Learning --- Bagging, Boosting, XGBOOST.
\normalsize

\vspace{0.3em}
\noindent\textbf{Unit V: Artificial Neural Networks}\\
\small
Artificial Neural networks --- Fundamental concepts on neural networks --- Activation Function --- Multilayer Feedforward Networks --- Convolutional Neural Network --- Recurrent Neural Network.
\normalsize

\vspace{0.5em}
\noindent\textbf{Books for Reference:}
\begin{enumerate}[leftmargin=1.5em, label=\arabic*., itemsep=1pt, topsep=1pt, parsep=0pt]
    \small
    \item Hao Ni, Xin Dong, Jinsons Zheng, Guangxi Yu (2021), \textit{An Introduction to Machine Learning in Quantitative Finance}, World Scientific.
    \item Mathew F Dixon, Igor Halperin, Paul Bilokon (2020), \textit{Machine Learning in Finance: From Theory to Practice}, Springer.
    \item Saurav Singla (2020), \textit{Machine Learning For Finance}, BPB Publications.
\end{enumerate}

\vspace{0.4em}
\begin{center}
\footnotesize
\begin{tabular}{|c|c|c|c|c|c|}
\hline
\textbf{COs} & \textbf{PO1} & \textbf{PO2} & \textbf{PO3} & \textbf{PO4} & \textbf{PO5} \\
\hline
\textbf{CO1} & X & XXX & XXX & XX & X \\
\hline
\textbf{CO2} & XX & XXX & XXX & XX & X \\
\hline
\textbf{CO3} & XX & XXX & XXX & XX & X \\
\hline
\textbf{CO4} & XX & XX & XXX & XXX & X \\
\hline
\textbf{CO5} & XX & XX & XX & XXX & X \\
\hline
\end{tabular}
\end{center}
\normalsize

\newpage
